% RESEARCH_DRAFT_SHA256: d7964dfaf84d972e00dcd3956a8a778ba4dc3d80f506f6e80dcc818c29234d04
% Layout and source-reference wrappers expanded for the web renderer.
\documentclass[11pt]{article}
\usepackage[margin=25mm]{geometry}
\usepackage{fontspec,unicode-math}
\setmainfont{Latin Modern Roman}
\setmathfont{Latin Modern Math}
\usepackage{amsmath,amssymb,mathtools,xurl,hyperref}
\hypersetup{colorlinks=true,urlcolor=blue,linkcolor=blue}
\setcounter{secnumdepth}{0}
\setlength{\parindent}{0pt}
\setlength{\parskip}{5pt}
\setlength{\emergencystretch}{3em}
\newcommand{\sref}[2]{\hyperlink{source:#1:#2}{[S#2]}}
\newcommand{\cataloguescope}[1]{\paragraph{Source record.}#1}
\newcommand{\floor}[1]{\left\lfloor#1\right\rfloor}
\begin{document}
\section*{Research-attempt provenance}
This standalone research attempt retains the complete original definition body below,
including its historical source and preparation notes. Those notes describe the
earlier catalogue preparation; the new attempt and its verification are identified
separately after that body. No original ID, catalogue status, or source statement is changed.

\section{Applications of Noncommutative F5 Gröbner Bases}\label{statement}

\subsection{Definitions and mathematical statement}

\paragraph{Shared notation.}Standard number systems are $\mathbb N,\mathbb Z,\mathbb Q,\mathbb R,\mathbb C$. Specialized notation and index ranges follow the cited source and the definitions below.


A Groebner basis generates an ideal and permits leading-term reductions under a chosen monomial order. In noncommutative settings word order and sidedness matter; an F5-type method uses signatures to avoid certain redundant reductions.
\paragraph{Record scope and qualifications.}
The catalogue title is retained as metadata. Named certificates, contracts, bounds or frameworks in that title are not asserted as results of the original source extract.
The record asks for potential applications, not a particular theorem or complexity guarantee.

Source-record limitation: Potential applications of an algorithm are requested without a mathematical proposition.
\paragraph{Statement recorded in the supplied source.}
What other mathematical applications are possible for a noncommutative F5 algorithm that computes Gröbner bases? \href{https://doi.org/10.4171/owr/2016/37}{[S2]}
\paragraph{Source status.}
The supplied catalogue labels this source record open; that metadata was checked against the supplied snapshot on 2026-10-11. The label is preserved for provenance and is not an independent verification of present open status or mathematical correctness. The catalogue literature assessment is dated 2026-08-21.
\subsection{Short English statement}
Find applications of the cited noncommutative Groebner-basis algorithm.
\subsection{Sources}
\begin{itemize}
\item[S1] Ulam AI and the UnsolvedMath Contributors, UnsolvedMath: A Curated Collection of Open Mathematics Problems, catalogue entry. Catalogue curation and metadata: CC BY 4.0. The source extract is retained; definitions, English summaries and qualifications are editorial additions.
\url{https://huggingface.co/datasets/ulamai/UnsolvedMath}.
\item[S2] Computational Group Theory (2017), Oberwolfach Reports, DOI 10.4171/owr/2016/37, pp. 2123--2169
\url{https://doi.org/10.4171/owr/2016/37}.
\item Source-specific Oberwolfach Reports licences, including CC BY-SA 4.0 where applicable, are retained; see THIRD\_PARTY\_NOTICES.md.
\end{itemize}
\paragraph{Preparation and review.}
The original source and UnsolvedMath contributors retain attribution for the source extract. Codex (GPT-6; exact serving revision not exposed) assisted with definitions, English summaries, scope notes and formatting. No new human mathematical authorship or independently verified solution is asserted. Catalogue identity checked against the complete supplied registry and current repository. Distinct upstream records, including repeated questions and broad or damaged extracts, are retained for complete coverage of the source CS category. A repeated record is not a new independent conjecture. Complete mathematical-equivalence and current-literature review has not been performed. Preparation and final review follow contribution prompts C and D, with source gaps disclosed. Structural and rendering checks do not establish a new theorem or open status.
\label{end:statement}

\section{Literature and context check}
Simon King's problem-session question in OWR 37/2016, pp.~2162--2163, was read. It says that a noncommutative F5 algorithm has been used for cohomology rings and asks for further applications. No signature rules, implementation interface, termination theorem, or benchmark is specified. The following worked quotient-algebra calculations describe one possible certificate-producing use of a noncommutative Groebner engine; they are standard elementary algebra rederived here, not evidence of a new F5 application or speedup.
\section{Research attempt: Finite quotient certificates and a noncommutative overlap test}

\paragraph{A complete small quotient calculation.}
Let $F$ be a field of characteristic different from two and work in the free associative algebra $F\langle x,y\rangle$ with a two-sided ideal
\[
 I=(x^2,y^2,yx+xy).
\]
Use reductions $xx\mapsto0$, $yy\mapsto0$, $yx\mapsto-xy$. They terminate: the first two lower word length, while the last preserves length and lowers the number of pairs in which a $y$ precedes an $x$. The only overlapping length-three ambiguities are $xxx,yyy,yyx,yxx$. The first two reduce to zero on either path; for example $yyx\mapsto-yxy\mapsto xyy\mapsto0$, and $yxx\mapsto-xyx\mapsto xxy\mapsto0$. Disjoint reductions commute, including their scalar signs.

These local checks imply a unique result: if two terminating reduction paths first choose different steps, disjoint steps or a checked overlap can be joined; induction on the maximum remaining reduction length then joins the rest. Every normal word is one of $1,x,y,xy$. Consequently every polynomial has a unique normal form
\[
 a+bx+cy+dxy,
\]
which vanishes if and only if it lies in $I$. One direction follows because every replacement subtracts a two-sided multiple of a defining relation. For the other direction, multiplication of a relation on either side reduces to zero by the just-proved uniqueness, so every element of the generated ideal has zero normal form. In particular the four normal words are linearly independent in the quotient.

This supplies explicit symbolic certificates. For $r=bx+cy+dxy$ the multiplication table gives $r^2=0$. Therefore $a+r$ is invertible exactly when $a\ne0$, with inverse $a^{-1}-a^{-2}r$. Necessity follows by mapping the quotient to $F$ via $x,y\mapsto0$. Also
\[
 [a+bx+cy+dxy,x]=-2cxy,\qquad
 [a+bx+cy+dxy,y]=2bxy.
\]
Thus the centralizer of $x$ is the span of $1,x,xy$, and the center is the span of $1,xy$. These are checkable applications to units and commutation constraints in a finitely presented associative algebra, provided the engine computes a terminating basis with this sidedness and order.

\paragraph{A concrete failed shortcut.}
It would be invalid to discard noncommutative overlaps merely because distinct leading words are not divisible by one another. For the relations $xy-x$ and $yx-y$, the word $xyx$ reduces by the left pair to $xx$ and by the right pair to $xy\mapsto x$. Thus the additional relation $x^2-x$ is forced; similarly $yxy$ forces $y^2-y$. With the four reductions
\[
 xx\mapsto x,\quad xy\mapsto x,\quad yx\mapsto y,\quad yy\mapsto y,
\]
any nonempty word reduces to its first letter. Every triple overlap has that same result and termination follows by length, so the repaired basis has normal words $1,x,y$. This exhibits an actual missing critical-pair calculation, not merely a warning that word order matters. A signature criterion may reject a pair only with a proof that its labelled consequence is already accounted for; the commutative ``no divisibility'' shortcut above is insufficient.
\section{Verification and adversarial audit}
All four overlap types in the first system and all eight triples in the repaired second system were checked symbolically by their first-letter rule. The characteristic restriction is needed for the center calculation, though not for the basic normal forms or nilpotence. No F5 implementation was run and no benchmark, generic termination claim, or performance comparison is asserted.
\paragraph{Audit corrections.} Made the two-sided free-associative setting explicit and supplied overlap/confluence arguments. Removed any implication that these known elementary quotient computations are new F5 performance results.
\paragraph{External computation.} None; symbolic arguments and boundary checks only.
\section{Scope and remaining gaps}
The original record is an open-ended applications question, not a theorem with a completion criterion. The worked finite-dimensional quotients show a concrete potential use and expose a bad pruning rule, but do not establish that the source's particular F5 variant supports this input or outperforms other methods. General presentations may lack finite bases, and their termination remains untreated.
\paragraph{Integration review.} Codex (GPT-6; exact serving revision not exposed) reviewed the written arguments and expanded layout and source-reference wrappers. Repository metadata and the production web renderer were checked using Python 3.13 and JavaScript with MathJax 3.2.2. This model-assisted review is not independent human proof verification; no PDF compilation or formal proof checking is asserted.

\section{Completion Estimate}
\noindent\textbf{COMPLETION: 0\%}
The calculations below the preserved source statement investigate a concrete route,
but establish no substantiated new progress toward the entire original target.
Known special cases and failed sufficient conditions are not credited as a solution.
\section{References}
The original source references and catalogue attribution are preserved in the
Sources section above. The following URLs were actually consulted for this attempt;
retention of an additional original reference does not claim fresh consultation of it.
\begin{itemize}
\item \url{https://ems.press/content/serial-article-files/46642} (consulted 2026-10-11). 
\end{itemize}
\section{Model and contributor disclosure}
OpenAI gpt-6-astra, requested reasoning configuration ultra, generated this research
attempt on 2026-10-11. The exact serving revision is not exposed. Codex GPT-6
performed orchestration, repository-format review, and written-proof review of this attempt; its exact serving revision
is also not exposed. The mathematical route, boundary checks and final adversarial
self-audit followed contribution prompts A then B in
\texttt{docs/contribution\_prompts.md}, after reading \texttt{CONTRIBUTING.md}.
Python 3.14.2 assembled the standalone documents; any mathematical computation is
specified in Verification. No theorem prover was used. No human mathematical
authorship, substantial human proof correction, or independent human proof checking
is asserted. Model-assisted review, source retrieval, and structural checks do not
establish correctness or novelty. Full private reasoning transcripts are unavailable
as reproducibility artifacts; the complete arguments, inputs and outcomes actually
used are stated here or in the companion batch audit.
\section{Final status}
Unresolved as a new research contribution. This is a concrete route investigation,
not a claim of new complete resolution or a current-status certification. A known
answer reproduced for context does not become a new contribution. The catalogue's
source-status label is preserved; this attempt status does not assert that a
source-labelled solved question is open.
\end{document}
